\chapter{Triiodothyronine (T3)}

\section*{What Is This Test?}
Triiodothyronine (T3) is the most biologically active thyroid hormone, with 3--5 times the metabolic potency of thyroxine (T4). Approximately 80\% of circulating T3 is produced by peripheral conversion (5'-deiodination) of T4 in the liver, kidney, and skeletal muscle by type 1 deiodinase (D1) and type 2 deiodinase (D2). Only about 20\% is directly secreted by the thyroid gland. In blood, T3 is extensively protein-bound (99.7\%): thyroid-binding globulin (TBG) carries approximately 65\%, transthyretin 10\%, and albumin 25\%. Only the unbound free fraction (FT3, approximately 0.3\%) is biologically active. The T3 test may measure total T3 (protein-bound plus free) or free T3, with free T3 preferred. Free T3 has a much shorter half-life (1 day) than free T4 (7 days). T3 is the primary ligand for nuclear thyroid hormone receptors (TR-alpha and TR-beta), regulating gene transcription for basal metabolic rate, cardiac contractility, thermogenesis, and neurological development. The test is most useful for confirming T3 thyrotoxicosis or monitoring anti-thyroid therapy.

\section*{Alternative Names}
Total T3; Total Triiodothyronine; Free T3 (FT3); Free Triiodothyronine; T3 Uptake; T3 Radioimmunoassay; Liothyronine (pharmaceutical T3)

\section*{Principle}
Free T3 is measured using a competitive immunoassay on automated platforms. The assay uses a labeled T3 analog (chemiluminescent tracer) that competes with free T3 in the patient sample for binding sites on an anti-T3 antibody. A displacement agent is used to prevent the labeled analog from binding to serum proteins. The chemiluminescent signal is inversely proportional to free T3 concentration. Common platforms: Roche Cobas ECLIA (electrochemiluminescence immunoassay, third-generation with improved low-end sensitivity), Abbott Architect CMIA, Siemens ADVIA Centaur/Atellica CLIA, and Beckman Coulter DxI. Total T3 is measured similarly but with a stronger displacement agent to release protein-bound T3. Equilibrium dialysis with LC-MS/MS detection is the gold standard reference method but is rarely used outside reference laboratories. Free T3 assays are technically challenging because T3 is present in very low concentrations (approximately 1/50th of T4 concentration), requiring high-sensitivity detection. Assay interference from T3 autoantibodies (rare) can produce spuriously elevated or decreased results.

\section*{History}
T3 was identified in 1952 by Jack Gross and Rosalind Pitt-Rivers during studies of thyroid hormone metabolism, demonstrating that T3 was more biologically active than T4. Jack Robbins and J. E. Rall discovered TBG binding of thyroid hormones in 1957. The first radioimmunoassay for T3 was developed by Chopra and colleagues in 1972. Free T3 assays using labeled antibody methods were developed in the 1980s. The clinical importance of T3 thyrotoxicosis (elevated T3 with normal T4 in hyperthyroidism) was recognized in the 1970s. In the 1990s, the role of peripheral T4-to-T3 conversion by deiodinase enzymes was elucidated by Reed Larsen and others. The importance of T3 measurement in non-thyroidal illness (low T3 syndrome) was described by Fliers and Boelen. Despite its biological importance, routine T3 measurement is not recommended for initial thyroid screening by ATA guidelines, and is reserved for specific clinical scenarios.

\section*{Why This Test Is Ordered}
\begin{itemize}
  \item Diagnosis of T3 thyrotoxicosis when TSH is suppressed but free T4 is normal --- T3 excess from Graves disease with disproportionate T3 secretion
  \item Evaluation of suspected toxic multinodular goiter or toxic adenoma where T3 may be elevated earlier than T4 in disease recurrence
  \item Monitoring response to anti-thyroid drug therapy (methimazole, propylthiouracil) in hyperthyroidism --- T3 normalizes before T4 during treatment
  \item Assessment of amiodarone-induced thyroid dysfunction --- amiodarone (Type 1 thyrotoxicosis) increases T3, while Type 2 (destructive thyroiditis) also releases preformed T3
  \item Investigation of suspected thyroid hormone resistance --- T3 may be disproportionately elevated
  \item Evaluation of T3-predominant autonomous thyroid function in patients with thyroid nodules
  \item Assessment of peripheral thyroid hormone metabolism in non-thyroidal illness (low T3 syndrome)
  \item Detection of surreptitious T3 (liothyronine, Cytomel) ingestion --- isolated T3 elevation with suppressed TSH and low T4/thyroglobulin
  \item Assessment of thyrotoxicosis in pregnancy when radioactive iodine uptake is contraindicated
  \item Serial monitoring during thyroid storm management (free T3 correlates with clinical severity)
\end{itemize}

\section*{Primary vs Secondary Test}
This is a secondary (adjunctive) test. Free T3 is not recommended as a screening test for thyroid dysfunction. The 2016 ATA hyperthyroidism guidelines and 2012 AACE/ATA hypothyroidism guidelines recommend TSH as the primary screening test, with free T4 as the first confirmatory test, and free T3 as an additional test in specific scenarios (elevated T4 with severe symptoms, suspected T3 thyrotoxicosis, monitored anti-thyroid drug therapy). Free T3 levels can be normal in up to 30\% of hypothyroid patients and are not useful for monitoring levothyroxine therapy.

\section*{How This Test Is Performed}
A healthcare professional draws a blood sample from a vein in your arm into a serum separator tube (gold-top) or plain red-top tube. Approximately 3--5 mL of blood is collected. The specimen is centrifuged to separate serum. Free T3 in serum is stable at 2--8 degrees C for 7 days and at -20 degrees C for 6 months. Avoid repeated freeze-thaw cycles. Specimens should be processed promptly; prolonged room-temperature storage can result in T3 degradation. Hemolyzed and icteric specimens should be avoided because hemoglobin and bilirubin interfere with chemiluminescent detection.

\section*{What Preparation Is Needed}
No fasting is required. Free T3 does not show clinically significant diurnal variation, so timing is less critical than for TSH. Patients on levothyroxine therapy should have blood drawn at a consistent time as part of their monitoring schedule. High-dose biotin (>5 mg/day) should be discontinued 48 hours before testing because it interferes with streptavidin-biotin assay platforms, potentially causing falsely elevated free T3 results. Patients receiving amiodarone should be monitored for thyroid dysfunction at baseline and every 6 months, including free T3 to detect amiodarone-induced thyrotoxicosis. The test should not be performed during acute illness unless non-thyroidal illness syndrome is specifically being investigated, since free T3 commonly decreases in hospitalized patients.

\section*{Contraindications and Precautions}
No absolute contraindications. Factors affecting T3 results: non-thyroidal illness/euthyroid sick syndrome (the commonest cause of low T3, due to decreased peripheral T4-to-T3 conversion by hepatic D1, increased reverse T3 production by D3); fasting and starvation (decrease T3 by 30--50\% within 24--48 hours as an adaptive response); severe hepatic disease (decreased D1 activity and reduced TBG synthesis); amiodarone (inhibits T4-to-T3 conversion, may decrease T3 while T4 is elevated); propylthiouracil (inhibits T4-to-T3 conversion more potently than methimazole); glucocorticoid therapy (suppresses TSH and inhibits T4-to-T3 conversion); beta-blockers (especially propranolol at high doses inhibit T4-to-T3 peripheral conversion); iodine-containing contrast agents (temporarily inhibit D1/D2, reducing T3); nephrotic syndrome (urinary loss of TBG lowers total T3 but not free T3); pregnancy (increased TBG from estrogen elevates total T3, while free T3 remains relatively stable). Heterophile antibodies and human anti-mouse antibodies can interfere with sandwich or competitive immunoassays, producing spurious T3 values. Free fatty acids from heparin therapy in vitro can displace T3 from binding proteins, artifactually elevating free T3 measurements --- if possible, blood should be drawn before heparin administration.

\section*{Risks}
Minimal risk. Slight pain, bruising, or hematoma at the venipuncture site. Rarely: vasovagal syncope or infection. No radiation exposure.

\section*{What Happens After the Test}
Results are available within 1--4 hours on routine automated immunoassay platforms. The healthcare provider interprets free T3 in conjunction with TSH, free T4, and clinical presentation. In hyperthyroid patients on anti-thyroid drugs, free T3 typically normalizes before TSH or free T4, so it is a useful early indicator of treatment efficacy. T3 thyrotoxicosis (TSH suppressed, FT4 normal, FT3 elevated) most commonly occurs in Graves disease and toxic nodular disease and requires initiation of anti-thyroid drugs (methimazole 10--30 mg/day). If T3-predominant toxicity is identified, beta-blockers (propranolol 10--40 mg every 6--8 hours) are indicated for symptomatic control due to its additional effect of inhibiting T4-to-T3 peripheral conversion.

\section*{Understanding Results}

\subsection*{Below Normal}
\begin{itemize}
  \item \textbf{Low free T3 with normal TSH and normal low free T4 (euthyroid sick syndrome/low T3 syndrome):} Most common cause of low T3 in hospitalized patients. Acute illness (sepsis, burns, trauma, major surgery, myocardial infarction, starvation) impairs T4-to-T3 conversion in liver via D1 downregulation. Reverse T3 is elevated. This is an adaptive response to conserve energy and protein and does not require thyroid hormone replacement. Low T3 correlates with poor prognosis in critically ill patients.
  \item \textbf{Low free T3 in severe hypothyroidism:} Overt primary hypothyroidism with decreased thyroidal T3 production and also decreased peripheral T4-to-T3 conversion. Free T3 may remain in the low-normal range in mild to moderate hypothyroidism because of compensatory upregulation of D2 activity in brain and pituitary.
  \item \textbf{Drug-induced low T3:} Amiodarone (potent D1 inhibitor), propylthiouracil (inhibits D1), glucocorticoids in high doses, propranolol at doses above 160 mg/day, iopanoic acid and ipodate (radiocontrast agents). Low T3 with normal/elevated T4 and normal TSH is expected.
  \item \textbf{Caloric deprivation and fasting:} T3 decreases within 24 hours as an energy conservation mechanism, with a corresponding increase in reverse T3.
  \item \textbf{Severe hepatic cirrhosis:} Decreased D1 activity, decreased TBG synthesis, and reduced thyroid hormone binding protein capacity.
  \item \textbf{Nephrotic syndrome:} Urinary loss of TBG and T3-binding proteins decreases total T3 (free T3 often normal).
\end{itemize}

\subsection*{Above Normal}
\begin{itemize}
  \item \textbf{T3 thyrotoxicosis:} Suppressed TSH, elevated free T3, and normal (or elevated) free T4. Accounts for approximately 5--10\% of hyperthyroid patients. Most commonly occurs in early Graves disease or toxic multinodular goiter when the thyroid gland produces T3 disproportionately due to increased D1/D2 expression in hyperfunctioning tissue. Also seen in iodine-deficient regions where T3 is preferentially synthesized.
  \item \textbf{Graves disease (overt hyperthyroidism):} Both free T4 and free T3 are elevated, but T3 may be disproportionately higher (T3:T4 ratio >20:1 in ng/ng units). TSH is fully suppressed (<0.01 mIU/L). Free T3 correlates well with symptoms of thyrotoxicosis including tremor, tachycardia, hyperreflexia, and heat intolerance.
  \item \textbf{Toxic adenoma (Plummer disease):} A single hyperfunctioning nodule autonomously producing T3 in excess of T4. TSH is suppressed, and the nodule is visible on thyroid ultrasound and radioactive iodine scan as a hot nodule with suppression of surrounding thyroid tissue.
  \item \textbf{Toxic multinodular goiter:} Multiple autonomous nodules producing mixed T4 and T3 secretion. T3 elevation may precede T4 elevation during disease evolution. More common in elderly patients with long-standing nodular goiter.
  \item \textbf{Exogenous T3 (liothyronine, Cytomel) ingestion:} Surreptitious or therapeutic T3 use causing elevated FT3, suppressed TSH, low or normal FT4, and low thyroglobulin (endogenous thyroid production suppressed). Unlike T4 ingestion, T3 elimination is rapid (half-life 1 day).
  \item \textbf{T3-producing metastatic thyroid carcinoma:} Rare functional metastases from follicular thyroid cancer can produce enough T3 to cause thyrotoxicosis.
  \item \textbf{Thyroid hormone resistance (TR-beta mutation):} Elevated free T3 and free T4 with inappropriately normal or elevated TSH. Patients may be asymptomatic or present with goiter, ADHD, tachycardia, or learning disability depending on tissue-specific resistance pattern.
  \item \textbf{TSH-secreting pituitary adenoma:} Elevated T3 and T4 with non-suppressed TSH. Alpha-subunit-to-TSH molar ratio is elevated.
  \item \textbf{Artifactual elevation from T3 autoantibodies:} Rare cases of T3 autoantibodies can interfere with immunoassays, producing falsely elevated total T3 results. Free T3 is less affected.
\end{itemize}

\section*{Normal Results}
\begin{itemize}
  \item \textbf{Free T3 (Adults, Roche ECLIA):} 2.0--4.4 pg/mL (3.1--6.8 pmol/L)
  \item \textbf{Free T3 (Adults, Abbott Architect):} 2.3--4.2 pg/mL (3.5--6.5 pmol/L)
  \item \textbf{Total T3 (Adults):} 80--200 ng/dL (1.2--3.1 nmol/L)
  \item \textbf{Total T3 (Pregnancy, third trimester):} 140--280 ng/dL (increased TBG)
  \item \textbf{Free T3 (Children 1--14 years):} 2.5--6.0 pg/mL
  \item \textbf{Free T3 (Neonates, 1--3 days):} 1.5--5.0 pg/mL (may be lower initially before TSH surge)
  \item \textbf{Free T3 (Elderly >70 years):} 1.8--3.8 pg/mL (lower due to decreased peripheral conversion with aging)
  \item \textbf{T3:T4 ratio (ng/ng units):} Normal <20:1; T3 thyrotoxicosis >20:1
  \item \textbf{Note:} Reference intervals vary significantly by assay methodology. Always use the laboratory-specific reference range provided with your results.
\end{itemize}

\section*{Follow-up Tests}
\begin{itemize}
  \item \textbf{TSH (third-generation):} Always co-interpreted with free T3. TSH alone is the best indicator of thyroid status --- suppressed TSH with elevated FT3 confirms thyrotoxicosis.
  \item \textbf{Free T4 (FT4):} Needed to distinguish T3 thyrotoxicosis (elevated FT3, normal FT4, suppressed TSH) from overt hyperthyroidism (both elevated) and to monitor treatment response.
  \item \textbf{Reverse T3 (rT3):} Evaluate low T3 syndrome vs. central hypothyroidism. In non-thyroidal illness, reverse T3 is elevated with low T3 and normal TSH. In central hypothyroidism, both T3 and reverse T3 are low.
  \item \textbf{TSH receptor antibodies (TRAb/TSI):} Confirm Graves disease as cause of T3-predominant thyrotoxicosis. Positive TRAb in >95\% of active Graves disease.
  \item \textbf{Radioactive iodine uptake (I-123) and scan:} Shows diffuse increased uptake (35--80\% at 24 hours) in Graves disease versus low/absent uptake in subacute thyroiditis or factitious T3 ingestion. A single hot nodule on scan indicates toxic adenoma.
  \item \textbf{Thyroid ultrasound with color Doppler:} Detect nodules in suspected toxic multinodular goiter. Increased parenchymal blood flow supports Graves disease.
  \item \textbf{Thyroglobulin:} Low serum thyroglobulin with elevated T3 and suppressed TSH suggests exogenous T3 ingestion rather than endogenous production.
  \item \textbf{Type 1 deiodinase (DIO1) and type 2 deiodinase (DIO2) genetic testing:} Research setting for thyroid hormone metabolism disorders.
  \item \textbf{Thyroid receptor beta (THRB) gene sequencing:} If thyroid hormone resistance suspected (elevated FT4 and FT3 with normal/elevated TSH).
\end{itemize}

\section*{References}
\begin{enumerate}
  \item Ross DS, Burch HB, Cooper DS, et al. 2016 American Thyroid Association guidelines for diagnosis and management of hyperthyroidism. Thyroid. 2016;26(10):1343-1421.
  \item Fliers E, Bianco AC, Langouche L, Boelen A. Thyroid function in critically ill patients. Lancet Diabetes Endocrinol. 2015;3(10):816-825.
  \item Brent GA. Mechanisms of thyroid hormone action. J Clin Invest. 2012;122(9):3035-3043.
  \item Bianco AC, Kim BW. Deiodinases: implications of the local control of thyroid hormone action. J Clin Invest. 2006;116(10):2571-2579.
  \item Jonklaas J, Bianco AC, Bauer AJ, et al. Guidelines for the treatment of hypothyroidism. Thyroid. 2014;24(12):1670-1751.
  \item Baloch Z, Carayon P, Conte-Devolx B, et al. Laboratory medicine practice guidelines for diagnosis and monitoring of thyroid disease. Thyroid. 2003;13(1):3-126.
  \item Gross J, Pitt-Rivers R. The identification of 3:5:3'-L-triiodothyronine in human plasma. Lancet. 1952;259(6705):439-441.
  \item Burtis CA, Ashwood ER, Bruns DE. Tietz Textbook of Clinical Chemistry and Molecular Diagnostics. 6th ed. Elsevier; 2023.
\end{enumerate}

\clearpage
\section*{Regulatory Status and Authorization Date}
Regulatory status applies to the specific assay, instrument, manufacturer, and intended use rather than to the test name alone. No single approval date can be assigned to this test category. For a commercial assay, verify the current product in the relevant regulator's database: the U.S. Food and Drug Administration (FDA) clearance, approval, or authorization; India's Central Drugs Standard Control Organisation (CDSCO) licence or permission; the European Union In Vitro Diagnostic Regulation (EU IVDR) CE certificate issued through the applicable conformity-assessment route; or the United Kingdom Medicines and Healthcare products Regulatory Agency (MHRA) registration and UKCA/CE status. Laboratory-developed tests may not have an individual FDA, CDSCO, EU IVDR, or MHRA approval date.
\clearpage
